\section{Exam benchmarks：难度可控但不一定真实}

考试类 benchmark 的吸引力很明显：题目边界清楚、答案通常唯一、评分便宜、难度可设计。老师把它类比成人类考试：考试可以控制学科和难度，但考试不等于真实工作。这个张力贯穿 MMLU、MMLU-Pro、GPQA 和 HLE。

本节的四组 benchmark 是一个递进序列：MMLU 建立了多学科知识考试范式；MMLU-Pro 解决饱和和噪声问题；GPQA 把题目推向专家级；HLE 则试图构造 frontier models 仍然很难的综合考试。读这些图时，要同时看两个维度：它们是否更难，以及它们是否更接近真实使用。后一维度并不会自动随难度增加。

\begin{figure}[H]
\centering
\includegraphics[width=0.88\textwidth]{mmlu.png}
\caption{MMLU：57 个 subject 的多选题，常被用作知识能力 benchmark。}
\figsource{CS336 Lecture 12 官方源引用图。}
\end{figure}

\begin{knowledgebox}{读图：MMLU 测的是知识，不是所有 language understanding}
MMLU 名字里有 language understanding，但其核心是跨学科知识问答。它的优点是题目清楚、容易评分、覆盖广；缺点是容易饱和、可能污染、且不代表真实助理交互。读这张图时，应把它理解为“知识考试面板”，而不是“通用智能温度计”。
\end{knowledgebox}

随着模型变强，旧考试会饱和。MMLU-Pro 的设计思路不是换一个名字，而是删掉噪声/简单题、增加选项数量、给模型 chain-of-thought 空间，让强模型仍有区分度。

\begin{figure}[H]
\centering
\includegraphics[width=0.88\textwidth]{mmlu-pro.png}
\caption{MMLU-Pro：去掉噪声/简单题，把 4 选项扩到 10 选项，并用 chain-of-thought 评估。}
\figsource{CS336 Lecture 12 官方源引用图。}
\end{figure}

\begin{knowledgebox}{读图：MMLU-Pro 为什么更难}
更多选项降低猜中概率，去掉 trivial questions 后减少饱和。模型准确率下降 16\%-33\% 说明旧 benchmark 的区分度已经不足。Benchmark 难度必须随模型进步而升级，否则 leaderboard 会变成“谁都差不多”的噪声表。
\end{knowledgebox}

GPQA 和 HLE 进一步把难度推高。GPQA 试图让非专家即使用 Google 也难以答对；HLE 则通过奖金、共同作者身份和多轮过滤构造极难题。它们的共同目标是恢复区分度，但代价是离普通用户使用越来越远。

\begin{figure}[H]
\centering
\includegraphics[width=0.78\textwidth]{gpqa.png}
\caption{GPQA：由 PhD contractors 编写的 graduate-level Google-proof Q\&A。}
\figsource{CS336 Lecture 12 官方源引用图。}
\end{figure}

\begin{knowledgebox}{读图：GPQA 的设计目标}
GPQA 试图让非专家即使用 Google 也难以答对，从而减少简单检索和记忆。它更接近专业知识推理，但题目来源、专家定义和评测格式仍会影响 validity。它测的是某种“专家题目上的知识和推理”，不是所有实际专业工作。
\end{knowledgebox}

\begin{figure}[H]
\centering
\includegraphics[width=0.78\textwidth]{hle-examples.png}
\caption{Humanity's Last Exam examples：多模态、多学科，包含多选和短答。}
\figsource{CS336 Lecture 12 官方源引用图。}
\end{figure}

\begin{figure}[H]
\centering
\includegraphics[width=0.9\textwidth]{hle-pipeline.png}
\caption{HLE pipeline：用奖金、共同作者身份、frontier LLM filtering 和多轮 review 构造高难题。}
\figsource{CS336 Lecture 12 官方源引用图。}
\end{figure}

\begin{figure}[H]
\centering
\includegraphics[width=0.72\textwidth]{hle-results.png}
\caption{HLE results：展示 frontier models 在更难考试题上的表现。}
\figsource{CS336 Lecture 12 官方源引用图。}
\end{figure}

\begin{knowledgebox}{读图：HLE 的意义和限制}
HLE 代表 benchmark 随模型变强而不断加难的趋势。它能延缓饱和，但仍然是考试式任务：答案通常预先存在，交互性弱，不一定覆盖真实世界开放式协作。老师的核心提醒是：更难不等于更真实；difficulty 和 realism 是两条不同轴。
\end{knowledgebox}

\subsection{本章小结}

Exam benchmarks 的优势是难度可控、答案明确、易评分；缺点是生态真实性弱，且容易被 saturation 和 contamination 追上。下一节转向 chat benchmarks：它们更真实，但评价更主观、更容易受偏好和 judge 影响。

